%%%%%%%%%%%%%%%%%%%%%%%%%%%%%%%%%%%%%%%%%
% Journal Article
% LaTeX Template
% Version 1.3 (9/9/13)
%
% This template has been downloaded from:
% http://www.LaTeXTemplates.com
%
% Original author:
% Frits Wenneker (http://www.howtotex.com)
%
% License:
% CC BY-NC-SA 3.0 (http://creativecommons.org/licenses/by-nc-sa/3.0/)
%
%%%%%%%%%%%%%%%%%%%%%%%%%%%%%%%%%%%%%%%%%

%----------------------------------------------------------------------------------------
%	PACKAGES AND OTHER DOCUMENT CONFIGURATIONS
%----------------------------------------------------------------------------------------

\documentclass[12pt]{article}

\usepackage{lipsum} % Package to generate dummy text throughout this template
\usepackage[french]{babel}

\usepackage[utf8]{inputenc}

%\usepackage[sc]{mathpazo} % Use the Palatino font
\usepackage[T1]{fontenc} % Use 8-bit encoding that has 256 glyphs
\linespread{1.05} % Line spacing - Palatino needs more space between lines
\usepackage{microtype} % Slightly tweak font spacing for aesthetics

\usepackage[hmarginratio=1:1,top=32mm,columnsep=20pt]{geometry} % Document margins
\usepackage{multicol} % Used for the two-column layout of the document
\usepackage[hang, small,labelfont=bf,up,textfont=it,up]{caption} % Custom captions under/above floats in tables or figures
\usepackage{booktabs} % Horizontal rules in tables
\usepackage{float} % Required for tables and figures in the multi-column environment - they need to be placed in specific locations with the [H] (e.g. \begin{table}[H])
\usepackage{hyperref} % For hyperlinks in the PDF
\usepackage{graphicx} 
\usepackage{lettrine} % The lettrine is the first enlarged letter at the beginning of the text
\usepackage{paralist} % Used for the compactitem environment which makes bullet points with less space between them

\usepackage{abstract} % Allows abstract customization
\renewcommand{\abstractnamefont}{\normalfont\bfseries}
\renewcommand{\abstractname}{Résumé} % Set the "Abstract" text to bold
\renewcommand{\abstracttextfont}{\normalfont\small\itshape} % Set the abstract itself to small italic text

\usepackage{titlesec} % Allows customization of titles
\renewcommand\thesection{\Roman{section}} % Roman numerals for the sections
\renewcommand\thesubsection{\arabic{subsection}.\arabic{subsection}} % Roman numerals for subsections
\titleformat{\section}[block]{\bfseries\large\centering}{\thesection.}{1em}{} % Change the look of the section titles
\titleformat{\subsection}[block]{\bfseries\large}{\thesubsection.}{1em}{} % Change the look of the section titles

\usepackage{fancyhdr} % Headers and footers
\pagestyle{fancy} % All pages have headers and footers
\fancyhead{} % Blank out the default header
\fancyfoot{} % Blank out the default footer
\renewcommand{\headrulewidth}{0pt} %pour enlever la ligne du header
%\fancyhead[C]{titre, date, noms...	} % Custom header text
%\fancyfoot[RO,LE]{\thepage} % Custom footer text


%agrandissement de la zone de texte
\addtolength{\oddsidemargin}{-1cm}
\addtolength{\evensidemargin}{-1cm}
\addtolength{\textwidth}{2cm}

%----------------------------------------------------------------------------------------
%	TITLE SECTION
%----------------------------------------------------------------------------------------

\title{ 
\includegraphics[width=60mm]{logo_JJCAB_2018.png}\\
\textbf{Débruitage de la matrice spectrale pour la caractérisation de sources aéroacoustiques}} % Article title

\author{
% KOMA Standard would be:
%  Peter Sciman\thanks{German Aerospace Center, Müller-Breslau-Str. 8, 10623 Berlin, Germany,
% peter.science@dlr.de}
% \and Sarah Techwoman\thanks{Technical University of Berlin, Straße des 17.~Juni 135,
% 10623 Berlin, Germany, sarah.technology@tu-berlin.de}
% Word template is better approximated by
 \normalsize{Alice Dinsenmeyer$^{1,2}$, Jérôme Antoni$^1$, Quentin Leclère$^1$}\\[-0.4em]
 \small{$^1$ Laboratoire Vibrations Acoustique}\\[-0.4em]
 \small{Univ Lyon, INSA-Lyon, LVA EA677, F-69621 Villeurbanne, France}\\[-0.4em]
 \small{$^2$ Laboratoire de Mécanique des Fluides et d’Acoustique }\\[-0.4em]
 \small{Univ Lyon, École Centrale de Lyon, F-69134, Écully, France}
\date{}
}
\date{}

%----------------------------------------------------------------------------------------

\begin{document}

\maketitle % Insert title

\thispagestyle{fancy} % All pages have headers and footers

%----------------------------------------------------------------------------------------
%	ABSTRACT
%----------------------------------------------------------------------------------------
\section*{Résumé}


Avec l'apparition des MEMS et la diminution globale du coût des capteurs, les acquisitions multivoies se généralisent, notamment dans le domaine de l'identification de sources acoustiques. La qualité de la localisation et de la quantification des sources peut être dégradée par la présence de bruit de mesure ambiant ou induit par le système d'acquisition. En particulier, dans le cas de mesures en présence d'un écoulement, le bruit de couche limite turbulente peut être supérieur au niveau des sources à caractériser, et il devient nécessaire de traiter les acquisitions pour extraire la contribution des sources de celle du bruit.

 Pour cela, on propose de décomposer la matrice spectrale mesurée en la somme d'une matrice signal et d'une matrice bruit. Cette décomposition exploite les propriétés statistiques de ces deux matrices. Comme le signal sources est corrélé sur les capteurs, le rang de la matrice spectrale associée se limite au nombre sources décorrélées équivalentes. De plus, le bruit aléatoire est supposé avoir une longueur de corrélation plus faible que celle du signal des sources, ce qui peut être approché par une matrice spectrale diagonale. En s'appuyant sur ce modèle, le débruitage est traité comme un problème d'optimisation bayésienne. 
 
À l'aide d'un cas numérique, les performances, en terme d'imagerie acoustique, sont comparées à la méthode classique de débruitage qui consiste à supprimer les éléments diagonaux de la matrice spectrale bruitée.
Dans un second temps, la méthode est appliquée à des données industrielles particulièrement bruitées, provenant de mesures microphoniques effectuées sur le fuselage d'un avion de ligne en vol.



\end{document}
